\documentclass[11pt]{amsart}
\usepackage{amsmath,amssymb,amsthm,booktabs}
\usepackage[margin=1.15in]{geometry}
\usepackage[hidelinks]{hyperref}
\newtheorem{theorem}{Theorem}[section]
\newtheorem{lemma}[theorem]{Lemma}
\newtheorem{conjecture}[theorem]{Conjecture}
\theoremstyle{remark}
\newtheorem{remark}[theorem]{Remark}
\newcommand{\R}{\mathbb{R}}
\newcommand{\HH}{\mathbb{H}}
\newcommand{\E}{\mathbb{E}}
\newcommand{\Prob}{\mathbb{P}}
\newcommand{\vol}{\operatorname{vol}}
\newcommand{\doi}[1]{\href{https://doi.org/#1}{doi:\nolinkurl{#1}}}
\newcommand{\arxiv}[1]{\href{https://arxiv.org/abs/#1}{arXiv:\nolinkurl{#1}}}
\setlength{\emergencystretch}{2.5em}

\title[Variance of the typical hyperbolic Poisson--Voronoi cell]{The variance of the typical cell of a hyperbolic
Poisson--Voronoi tessellation at low intensity}
\author{Yitzchak Shmalo}
\address{Yitzchak Shmalo, Tel Aviv University and Heights Labs}
\email{yspennstate@gmail.com}
\date{November 2026}
\subjclass[2020]{60D05, 60G55, 52A55}
\keywords{Poisson--Voronoi tessellation, hyperbolic space, ideal Poisson--Voronoi tessellation, typical cell, zero
cell, size biasing}

\raggedbottom
\begin{document}
\begin{abstract}
Let $V_d(\lambda)$ be the variance of the volume of the typical cell of the Poisson--Voronoi tessellation of the
hyperbolic space $\HH^d$ with intensity $\lambda$. D'Achille, Curien, Enriquez, Lyons and \"Unel asked whether
$\lambda V_d(\lambda)$ has a finite limit as $\lambda\to0$, equal to the expected volume of
$C_d\setminus T_d(R_1)$, where $C_d$ is the zero cell of the ideal Poisson--Voronoi tessellation and $T_d(R_1)$ the
set coupled with it. We show that $\lambda V_d(\lambda)\to\infty$ for every $d\ge2$ and that the expected volume of
$C_d\setminus T_d(R_1)$ is infinite, growing like $\theta_d\log U$ in the window $\{|u|<U\}$ of the upper half-space
model, with $\theta_2=8/\pi$. The proof couples the typical cell inside the zero cell at positive intensity and passes
to the limit with Fatou's lemma. In the plane, numerical values of $\lambda V_2(\lambda)$ down to $\lambda=10^{-5}$ fit
$(8/\pi)\log(1/\lambda)-6.84$.
\end{abstract}
\maketitle

\section{Introduction}\label{sec:intro}

Let $X^{(\lambda)}$ be a Poisson point process of intensity $\lambda\vol$ on the hyperbolic space $\HH^d$ of constant
curvature $-1$, $d\ge2$, and let $\mathrm{Vor}(X^{(\lambda)})$ be its Voronoi tessellation. Unlike the Euclidean case,
these tessellations have a nontrivial limit in law as $\lambda\to0$, the ideal Poisson--Voronoi tessellation of
D'Achille, Curien, Enriquez, Lyons and \"Unel~\cite{DCELU}. Low-intensity Poisson--Voronoi tessellations of hyperbolic
surfaces were used by Budzinski, Curien and Petri~\cite{BCP} to bound Cheeger constants, and the ideal tessellation of
higher-rank symmetric spaces by Fraczyk, Mellick and Wilkens~\cite{FMW} to prove fixed price. Mellick~\cite{Mellick}
showed that the cells of the ideal tessellation on $\HH^d$ are indistinguishable, D'Achille~\cite{DA} constructed ideal
tessellations beyond hyperbolic spaces, and D'Achille and Th\"ale~\cite{DT} computed face volume densities at positive
and zero intensity; none of these concerns the variance of the typical cell. In dimension three Isokawa~\cite{Isokawa}
computed mean values for the cells at positive intensity.

The zero cell $C_d$ of the ideal tessellation, the cell containing a fixed origin $o$, has an explicit description~\cite[Theorem~1.4]{DCELU}: in the upper
half-space model with the unique end of $C_d$ sent to $\infty$, $C_d$ is the complement of a Poisson cloud of open
half-balls centred on the boundary, conditioned on not covering $o$ (Section~\ref{sec:setting}). Removing the
condition gives a stationary random set $T_d(R_1)$ which describes $C_d$ far from $o$ and is contained in it.

Let $V$ be the volume of the typical cell of $\mathrm{Vor}(X^{(\lambda)})$, so that $\E V=1/\lambda$, and let
$V_d(\lambda)=\operatorname{Var}V$. The zero cell has the size-biased law of the typical cell, so its expected volume is
$\E V^2/\E V=\lambda V_d(\lambda)+1/\lambda$, and $\lambda V_d(\lambda)$ is the difference between the expected volumes
of the zero cell and of the typical cell. This led to the following question.

\medskip\noindent\textbf{Question}~\cite[Question~7.7]{DCELU}. Does $\lambda V_d(\lambda)$ have a finite limit as
$\lambda\to0$? Is the limit equal to the expected volume of $C_d\setminus T_d(R_1)$ in their natural coupling? If so,
what is this value?
\medskip

\begin{theorem}\label{thm:main}
For every $d\ge2$, $\lambda V_d(\lambda)\to\infty$ as $\lambda\to0$.
\end{theorem}

\begin{theorem}\label{thm:coupling}
For every $d\ge2$, $\E\vol\bigl(C_d\setminus T_d(R_1)\bigr)=\infty$. More precisely, with $W_U=\{(u,v)\in\R^{d-1}\times
(0,\infty):|u|<U\}$ in the upper half-space model of Section~\ref{sec:setting},
\[
\lim_{U\to\infty}\frac{\E\vol\bigl((C_d\setminus T_d(R_1))\cap W_U\bigr)}{\log U}=\theta_d:=
\frac{2^{d+1}\pi^{(d-2)/2}\,\Gamma(d/2)}{(d-1)^2\,\Gamma\bigl(\tfrac{d-1}2\bigr)^2},
\]
so that $\theta_2=8/\pi$, $\theta_3=2\pi$ and $\theta_4=128/9$.
\end{theorem}

Theorem~\ref{thm:main} answers the first part of the question negatively, and Theorem~\ref{thm:coupling} shows that
the proposed value is infinite. The heuristic behind the question survives in the following form. At positive
intensity there is an exact coupling in which the typical cell is a subset of the zero cell and $\lambda V_d(\lambda)$
is the expected volume of the difference (Lemma~\ref{lem:coupling}); as $\lambda\to0$ the probability that a given
point lies in this difference converges to the corresponding probability for the natural coupling of $C_d$ and
$T_d(R_1)$ (Lemma~\ref{lem:limits}), and Fatou's lemma gives
\[
\liminf_{\lambda\to0}\lambda V_d(\lambda)\ \ge\ \E\vol\bigl(C_d\setminus T_d(R_1)\bigr)=\infty.
\]
The divergence of the coupling volume comes from points far from the vertical axis through $o$. A point $w=(u,v)$ lies
in $C_d\setminus T_d(R_1)$ only if it is covered by a half-ball containing $o$ and by none that misses $o$, and for $|u|$ large the
integral of this probability over $v$ is of order $|u|^{1-d}$; summed over the window $W_U$ this gives $\theta_d\log U$.

In the plane we computed $\lambda V_2(\lambda)$ numerically from an exact integral formula, with an independent Monte
Carlo check (Section~\ref{sec:numerics}). The values increase by $(8/\pi)\log10$ per decade of $1/\lambda$ in the
limit, the constant of Theorem~\ref{thm:coupling}.

\begin{conjecture}\label{conj:rate}
$\lambda V_2(\lambda)=\frac8\pi\log\frac1\lambda+C+o(1)$ as $\lambda\to0$, with $C\approx-6.84$.
\end{conjecture}

The two growth rates agree if the finite-intensity coupling matches the ideal one inside a window $W_U$ with $U$ of
order $1/\lambda$ and the contribution from outside the window stays bounded. We have no numerics in dimension $d\ge3$
and do not know the rate there.

Some conventions. Part of~\cite{DCELU} parametrizes the intensity as $\lambda^{d-1}\vol$; Theorem~\ref{thm:main}
holds with either normalization, since with intensity $\mu=\lambda^{d-1}$ one has $\lambda V_d(\mu)=\mu^{(2-d)/(d-1)}\,(\mu
V_d(\mu))$ and $\mu^{(2-d)/(d-1)}\ge1$ for $\mu\le1$. Throughout, $|S^{k}|=2\pi^{(k+1)/2}/\Gamma(\frac{k+1}2)$ is the area of the
unit sphere of $\R^{k+1}$, with $|S^0|=2$.

Section~\ref{sec:setting} recalls the description of $C_d$ and writes the expected volume of $C_d\setminus T_d(R_1)$ as
an integral. Section~\ref{sec:finite} gives the coupling at positive intensity, Section~\ref{sec:limit} its limit and
the proof of Theorem~\ref{thm:main}, and Section~\ref{sec:coupling} the proof of Theorem~\ref{thm:coupling}.
Section~\ref{sec:numerics} reports the planar computations, and Appendix~\ref{app:notation} lists the notation.

\section{The ideal zero cell and its stationary companion}\label{sec:setting}

We use the upper half-space model $\mathcal U_d=\R^{d-1}\times(0,\infty)$, points $w=(u,v)$, volume element
$du\,dv/v^d$ and origin $o=(0,1)$. For $x\in\R^{d-1}$ and $\rho>0$ let $B(x,\rho)=\{w\in\mathcal U_d:|w-x|<\rho\}$
be the open Euclidean half-ball centred at the boundary point $(x,0)$; here and below $|w-x|^2=|u-x|^2+v^2$. On
$\R^{d-1}\times(0,\infty)$ let
\[
m(dx\,d\rho)=\rho^{1-2d}\,dx\,d\rho,\qquad A_w=\{(x,\rho):w\in B(x,\rho)\}=\{(x,\rho):\rho>|w-x|\}.
\]
Let
\[
c_d=\frac{2^{2-d}\pi^{d/2}}{\Gamma(d/2)}=2^{1-d}|S^{d-1}|,\qquad M=\frac{c_d}{2(d-1)}.
\]

\begin{lemma}\label{lem:measures}
$m(A_w)=Mv^{1-d}$ for every $w=(u,v)$; in particular $m(A_o)=M$.
\end{lemma}

\begin{proof}
Integrating out $\rho$, $m(A_w)=\frac1{2d-2}\int_{\R^{d-1}}(|u-x|^2+v^2)^{1-d}\,dx$, and the substitution $x=u+vy$ gives
$v^{1-d}\frac1{2d-2}\int_{\R^{d-1}}(1+|y|^2)^{1-d}dy$. The last integral is $\pi^{(d-1)/2}\Gamma(\frac{d-1}2)/\Gamma(d-1)$,
which equals $c_d$ by the duplication formula $\Gamma(d-1)=2^{d-2}\pi^{-1/2}\Gamma(\frac{d-1}2)\Gamma(\frac d2)$.
\end{proof}

By~\cite[Theorem~1.4 and its proof in Section~5.1]{DCELU}, $C_d$ can be realised as follows. Let $R_1$ be an
$\mathrm{Exp}(1)$ random variable and, given $R_1=s$, let $N$ be a Poisson process on $\R^{d-1}\times(0,\infty)$ with
intensity $\kappa\,m$, where
\[
\kappa=\frac{2(d-1)s}{c_d}=\frac sM .
\]
Then $C_d$ has the law of the complement of $\bigcup\{B(x,\rho):(x,\rho)\in N,\ \rho^2\le1+|x|^2\}$, the union of the
half-balls of $N$ that do not contain $o$. In~\cite[Section~5.3]{DCELU}, $T_d(s)$ is the complement of the union of all
the half-balls of $N$; in the natural coupling $C_d$ and $T_d(R_1)$ are built from the same $N$, so $T_d(R_1)\subset
C_d$, and $T_d(R_1)$ uses the $\mathrm{Poisson}(\kappa m(A_o))=\mathrm{Poisson}(R_1)$ further half-balls that contain $o$.
(The quantity $\kappa\rho^{1-2d}$ is the intensity $2(d-1)s\,c_d^{-1}\rho^{1-2d}$ of~\cite[(5.4)]{DCELU}; their $\varrho$
is our $\rho$.)

For $w\in\mathcal U_d$ put $I(w)=m(A_w\cap A_o)$. Since the boundaries of $A_w$ and $A_o$ are $m$-null,
\begin{align*}
\Prob\bigl(w\in C_d\setminus T_d(R_1)\bigm|R_1=s\bigr)&=\Prob\bigl(N(A_w\setminus A_o)=0\bigr)-\Prob\bigl(N(A_w)=0\bigr)\\
&=e^{-\kappa(m(A_w)-I(w))}-e^{-\kappa m(A_w)} .
\end{align*}
With $\kappa m(A_w)=sv^{1-d}$, $\kappa I(w)=sI(w)/M$ and $\E e^{-aR_1}=1/(1+a)$, Fubini's theorem gives
\begin{equation}\label{eq:coupling-integral}
\E\vol\bigl((C_d\setminus T_d(R_1))\cap W\bigr)=\int_W\Bigl[\frac1{1+v^{1-d}-I(w)/M}-\frac1{1+v^{1-d}}\Bigr]
\frac{du\,dv}{v^d}
\end{equation}
for every Borel set $W\subset\mathcal U_d$. The bracket is nonnegative because $0\le I(w)\le m(A_w)=Mv^{1-d}$.

\section{A coupling at positive intensity}\label{sec:finite}

Fix $\lambda>0$ and let $X=X^{(\lambda)}$. Let $X_1$ be the point of $X$ closest to $o$, $r=d(o,X_1)$, and $B_o=
B_{\HH}(o,r)$ the open hyperbolic ball. Write $\mathcal V(\varrho)=|S^{d-1}|\int_0^\varrho\sinh^{d-1}t\,dt$ for the
volume of a hyperbolic ball of radius $\varrho$.

\begin{lemma}\label{lem:coupling}
Let $\Pi_{\rm in}$ be, given $X$, a Poisson process of intensity $\lambda\vol$ on $B_o$, and $\Pi=(X\setminus\{X_1\})\cup
\Pi_{\rm in}$. Let $C^\lambda$ be the cell of $X_1$ in $\mathrm{Vor}(X)$, that is the zero cell, and $T^\lambda$ the
cell of $X_1$ in $\mathrm{Vor}(\{X_1\}\cup\Pi)$. Then
\begin{enumerate}
\item $T^\lambda\subset C^\lambda$;
\item $\vol(T^\lambda)$ has the law of the volume of the typical cell of $\mathrm{Vor}(X^{(\lambda)})$;
\item $\lambda V_d(\lambda)=\E\vol(C^\lambda\setminus T^\lambda)$;
\item with $B_z=B_{\HH}(z,d(z,X_1))$,
\[
\Prob\bigl(z\in C^\lambda\setminus T^\lambda\bigm|X_1\bigr)=e^{-\lambda\vol(B_z\setminus B_o)}-e^{-\lambda\vol(B_z)}
\quad\text{for every }z\in\HH^d .
\]
\end{enumerate}
\end{lemma}

\begin{proof}
(1) Adding nuclei can only shrink the cell of $X_1$. (2) Given $X_1$, the process $X\setminus\{X_1\}$ is a Poisson
process of intensity $\lambda\vol$ on $\HH^d\setminus\overline{B_o}$, and $\Pi_{\rm in}$ is an independent Poisson
process of the same intensity on $B_o$. Hence, given $X_1$, $\Pi$ is a Poisson process of intensity $\lambda\vol$ on
$\HH^d$; its conditional law does not depend on $X_1$, so $\Pi$ is independent of $X_1$. By the Mecke equation~\cite[Theorem~4.1]{LP}
(Slivnyak's theorem), the typical cell is the cell of a fixed point $y$ in $\mathrm{Vor}(\{y\}\cup\Pi)$,
and since the isometry group acts transitively on $\HH^d$ its volume has the same law for every $y$. Conditioning on
$X_1$ gives (2). (3) For $x\in\HH^d$ let $\iota_x$ be the geodesic symmetry about the midpoint of $o$ and $x$, an
isometry exchanging $o$ and $x$, and write $C_x(Y)$ for the cell of $x$ in $\mathrm{Vor}(Y)$. By the Mecke equation and
the invariance of the law of $X$ under $\iota_x$,
\begin{align*}
\E\vol(C^\lambda)&=\E\sum_{x\in X}\vol(C_x(X))\mathbf 1\{o\in C_x(X)\}\\
&=\lambda\int_{\HH^d}\E\bigl[\vol(C_o(X\cup\{o\}))\mathbf 1\{x\in C_o(X\cup\{o\})\}\bigr]dx=\lambda\,\E V^2,
\end{align*}
so the zero cell has the size-biased law of the typical cell, as in Euclidean space~\cite[Section~10.4]{SW}. Hence
$\E\vol(C^\lambda)=\lambda V_d(\lambda)+1/\lambda$ in $[0,\infty]$, while $\E\vol(T^\lambda)=1/\lambda$ by (2), and (1)
gives (3). (4) Up to null events, $z\in C^\lambda$ if and only if $X\setminus\{X_1\}$
has no point in $B_z$, and $z\in T^\lambda$ if and only if $\Pi$ has no point in $B_z$. Given $X_1$ these events have
probabilities $e^{-\lambda\vol(B_z\setminus B_o)}$ and $e^{-\lambda\vol(B_z)}$, and the second event is contained in the
first.
\end{proof}

By (3), (4) and Tonelli's theorem,
\[
\lambda V_d(\lambda)=\int_{\HH^d}\E\bigl[e^{-\lambda\vol(B_z\setminus B_o)}-e^{-\lambda\vol(B_z)}\bigr]\,dz .
\]
The integrand depends on $X_1$ only through $r$ and the direction of $X_1$ seen from $o$; the direction is uniform and
independent of $r$, and $S=\lambda\mathcal V(r)$ is $\mathrm{Exp}(1)$. By rotation invariance about $o$ we may fix the
direction. Identify $\HH^d$ with $\mathcal U_d$ by an isometry that sends $o$ to $(0,1)$ and the ideal endpoint of that
direction to $\infty$, so that $X_1=\xi_s:=(0,e^{r})$ with $\lambda\mathcal V(r)=s$. Writing $B_w$, $B_o$ for the balls
above with $X_1=\xi_s$,
\begin{equation}\label{eq:finite-integral}
\lambda V_d(\lambda)=\int_0^\infty e^{-s}\,ds\int_{\mathcal U_d}g_\lambda(w,s)\,\frac{du\,dv}{v^d},\qquad
g_\lambda(w,s)=e^{-\lambda\vol(B_w\setminus B_o)}-e^{-\lambda\vol(B_w)}\ \ge0 .
\end{equation}
Here $r=r(\lambda,s)\to\infty$ as $\lambda\to0$ for fixed $s$.

\section{Passage to the limit}\label{sec:limit}

We use two facts of hyperbolic geometry. First, $\mathcal V(\varrho)e^{-(d-1)\varrho}\to|S^{d-1}|2^{1-d}/(d-1)$ as
$\varrho\to\infty$, so that
\begin{equation}\label{eq:ratio}
\frac{\mathcal V(\varrho)}{\mathcal V(\varrho')}\to e^{(d-1)\delta}\qquad\text{whenever }\varrho,\varrho'\to\infty,\
\varrho-\varrho'\to\delta .
\end{equation}
Second, for a boundary point $\eta$ of $\mathcal U_d$ and a geodesic ray $t\mapsto c(t)$ from $o$ to $\eta$, the
Busemann function $b_\eta(w)=\lim_{t\to\infty}\bigl(d(w,c(t))-t\bigr)$ exists, and in our coordinates
\begin{equation}\label{eq:busemann}
b_\infty(w)=-\log v,\qquad b_x(w)=\log\frac{|w-x|^2}{v\,(1+|x|^2)}\quad(x\in\R^{d-1}) .
\end{equation}
The first follows from $d((u,v),(0,e^t))=\operatorname{arccosh}\bigl(1+\frac{|u|^2+(v-e^t)^2}{2ve^t}\bigr)=t-\log
v+o(1)$; the second follows from the first and an isometry of $\mathcal U_d$ fixing $o$ and sending $\infty$ to $x$,
since horospheres centred at $x$ are the Euclidean spheres tangent to the boundary at $x$ (see~\cite[Section~4.6]{Rat}).
The directions at $o$ correspond to boundary points: if $\omega$ is uniform on the unit sphere of $T_o\HH^d$, the
endpoint $x(\omega)$ of the ray from $o$ in direction $\omega$ has density $c_d^{-1}(1+|x|^2)^{1-d}$ on $\R^{d-1}$
\cite[Lemma~3.1]{DCELU}.

\begin{lemma}\label{lem:limits}
Fix $s>0$ and $w=(u,v)\in\mathcal U_d$. As $\lambda\to0$,
\[
\lambda\vol(B_w)\to sv^{1-d},\qquad\lambda\vol(B_w\cap B_o)\to\frac{s\,I(w)}M .
\]
Consequently $g_\lambda(w,s)\to e^{-s(v^{1-d}-I(w)/M)}-e^{-sv^{1-d}}$.
\end{lemma}

\begin{proof}
Since $\lambda\mathcal V(r)=s$, $\lambda\vol(B_w)=s\,\mathcal V(d(w,\xi_s))/\mathcal V(r)$, and $d(w,\xi_s)-r\to
b_\infty(w)=-\log v$ by~\eqref{eq:busemann}; \eqref{eq:ratio} gives the first limit.

For the second, write a point of $B_o$ in geodesic polar coordinates about $o$ as $y=\exp_o(\varrho\omega)$ and put
$t=\lambda\mathcal V(\varrho)\in(0,s)$. Then $\lambda\,dy=dt\,\sigma(d\omega)$ with $\sigma$ the uniform probability on
directions, and
\[
\lambda\vol(B_w\cap B_o)=\int_0^s dt\int\sigma(d\omega)\,\mathbf 1\{d(w,y)<d(w,\xi_s)\} .
\]
For fixed $\omega$ and $t\in(0,s)$, $\varrho=\mathcal V^{-1}(t/\lambda)\to\infty$ and $\varrho-r\to\frac1{d-1}\log
\frac ts$ by~\eqref{eq:ratio}, so
\[
d(w,y)-d(w,\xi_s)=\bigl(d(w,y)-\varrho\bigr)-\bigl(d(w,\xi_s)-r\bigr)+(\varrho-r)\ \to\ b_{x(\omega)}(w)+\log
v+\frac1{d-1}\log\frac ts .
\]
For fixed $\omega$ the limit is strictly increasing in $t$, so it vanishes on a null set of $(t,\omega)$, and the
indicators converge almost everywhere on the finite measure space $(0,s)\times S^{d-1}$. By dominated convergence and
\eqref{eq:busemann},
\[
\lambda\vol(B_w\cap B_o)\to\int_0^s dt\int_{\R^{d-1}}\frac{dx}{c_d(1+|x|^2)^{d-1}}\,
\mathbf 1\Bigl\{|w-x|^2<(1+|x|^2)\Bigl(\frac st\Bigr)^{1/(d-1)}\Bigr\} .
\]
Substitute $\rho=(1+|x|^2)^{1/2}(s/t)^{1/(2d-2)}$, which maps $t\in(0,s)$ onto $\rho>(1+|x|^2)^{1/2}$ with
$dt=-2(d-1)s(1+|x|^2)^{d-1}\rho^{1-2d}d\rho$; the minus sign reverses the limits of integration. The limit becomes
\[
\frac{2(d-1)s}{c_d}\int_{\R^{d-1}}dx\int_{(1+|x|^2)^{1/2}}^\infty\mathbf 1\{\rho>|w-x|\}\,\rho^{1-2d}\,d\rho
=\kappa\,m(A_w\cap A_o)=\frac{sI(w)}M .
\]
Finally $\vol(B_w\setminus B_o)=\vol(B_w)-\vol(B_w\cap B_o)$.
\end{proof}

\begin{proof}[Proof of Theorem~\ref{thm:main}]
The integrand in~\eqref{eq:finite-integral} is nonnegative, so Fatou's lemma and Lemma~\ref{lem:limits} give
\[
\liminf_{\lambda\to0}\lambda V_d(\lambda)\ge\int_0^\infty e^{-s}ds\int_{\mathcal U_d}\Bigl(e^{-s(v^{1-d}-I(w)/M)}-
e^{-sv^{1-d}}\Bigr)\frac{du\,dv}{v^d}=\E\vol\bigl(C_d\setminus T_d(R_1)\bigr),
\]
the last equality by~\eqref{eq:coupling-integral} with $W=\mathcal U_d$. The right side is infinite by
Theorem~\ref{thm:coupling}.
\end{proof}

\section{The coupling volume}\label{sec:coupling}

By rotation invariance about the vertical axis, $I(u,v)$ depends on $u$ only through $|u|$. Write
\[
J(u)=\int_0^\infty\Bigl[\frac1{a-\varepsilon}-\frac1a\Bigr]\frac{dv}{v^d},\qquad a=1+v^{1-d},\quad
\varepsilon=\varepsilon(u,v)=\frac{I(u,v)}M\in[0,v^{1-d}],
\]
so that by~\eqref{eq:coupling-integral} $\E\vol((C_d\setminus T_d(R_1))\cap W_U)=\int_{|u|<U}J(u)\,du$. Since
$A_w\cap A_o=\{(x,\rho):\rho>\max(|w-x|,(1+|x|^2)^{1/2})\}$,
\begin{equation}\label{eq:I}
I(u,v)=\frac1{2d-2}\int_{\R^{d-1}}\max\bigl(|u-x|^2+v^2,\ |x|^2+1\bigr)^{1-d}dx .
\end{equation}

\begin{lemma}\label{lem:I}
\begin{enumerate}
\item $J(u)\le1+\frac1{2d-2}$ for every $u$.
\item $|u|^{d-1}I(u,v)\le K_d:=2^{d-1}|S^{d-2}|/(d-1)^2$ for every $u\ne0$ and $v>0$.
\item For every $v>0$, $|u|^{d-1}I(u,v)\to L_d$ as $|u|\to\infty$, where
\[
L_d=\frac1{2d-2}\int_{\R^{d-1}}\max(|y-e|^2,|y|^2)^{1-d}dy=\frac{2^{d-1}\pi^{(d-2)/2}\,\Gamma(d/2)}{(d-1)^2\,\Gamma(d-1)}
\]
for any unit vector $e$.
\end{enumerate}
\end{lemma}

\begin{proof}
(1) We have $\varepsilon\le\min(1,v^{1-d})$, because $I\le m(A_o)=M$ and $I\le m(A_w)=Mv^{1-d}$. So $a-\varepsilon\ge
\max(1,v^{1-d})$ and the bracket is $\varepsilon/(a(a-\varepsilon))$. For $v\le1$ it is at most $1/(a\,v^{1-d})\le
v^{2d-2}$, and $v^{2d-2}v^{-d}\le1$. For $v>1$ it is at most $\varepsilon/a\le v^{1-d}$, and $\int_1^\infty v^{1-2d}dv=
\frac1{2d-2}$.

(2) Split $\R^{d-1}$ into $H_1=\{|x|\le|x-u|\}$ and $H_2=\{|x|>|x-u|\}$. On $H_1$, $|x-u|\ge|u|/2$ and the max in
\eqref{eq:I} is at least $|x-u|^2$; on $H_2$, $|x|\ge|u|/2$ and it is at least $|x|^2$. Each part is at most
$\int_{|y|\ge|u|/2}|y|^{2-2d}dy=|S^{d-2}|(|u|/2)^{1-d}/(d-1)$.

(3) Substituting $x=|u|y$ and $u=|u|e$ in~\eqref{eq:I},
\[
|u|^{d-1}I(u,v)=\frac1{2d-2}\int_{\R^{d-1}}\max\Bigl(|y-e|^2+\frac{v^2}{|u|^2},\ |y|^2+\frac1{|u|^2}\Bigr)^{1-d}dy,
\]
which increases to $L_d$ by monotone convergence. On the half-space $\{y\cdot e\le\frac12\}$ the max is $|y-e|^2$,
on the other half it is $|y|^2$, and the two halves are exchanged by $y\mapsto e-y$; so $L_d=\frac1{d-1}\int_{z\cdot e
\ge1/2}|z|^{2-2d}dz$. Writing $z=(z_1,z')$ with $z_1=z\cdot e$ and $z'=z_1\zeta$,
\[
\int_{z_1\ge1/2}|z|^{2-2d}dz=\int_{1/2}^\infty z_1^{-d}dz_1\int_{\R^{d-2}}(1+|\zeta|^2)^{1-d}d\zeta
=\frac{2^{d-1}}{d-1}\cdot\frac{\pi^{(d-2)/2}\Gamma(d/2)}{\Gamma(d-1)},
\]
where for $d=2$ the inner integral is $1$.
\end{proof}

\begin{proof}[Proof of Theorem~\ref{thm:coupling}]
Fix $u$ with $|u|$ so large that $K_d|u|^{1-d}\le M/2$. Then $\varepsilon\le\frac12\le\frac a2$ by
Lemma~\ref{lem:I}(2), so $a-\varepsilon\ge a/2$ and
\[
|u|^{d-1}\Bigl[\frac1{a-\varepsilon}-\frac1a\Bigr]v^{-d}=\frac{|u|^{d-1}\varepsilon}{a(a-\varepsilon)v^d}\le
\frac{2K_d}M\,\frac{v^{-d}}{a^2}=\frac{2K_d}M\,\frac{v^{d-2}}{(1+v^{d-1})^2},
\]
which is integrable on $(0,\infty)$. As $|u|\to\infty$, $\varepsilon\to0$ and $|u|^{d-1}\varepsilon\to L_d/M$ for
every $v$ by Lemma~\ref{lem:I}(3), so by dominated convergence
\[
|u|^{d-1}J(u)\to\frac{L_d}M\int_0^\infty\frac{v^{d-2}\,dv}{(1+v^{d-1})^2}=\frac{L_d}{(d-1)M}.
\]
Since $J$ is bounded (Lemma~\ref{lem:I}(1)) and radial,
\[
\int_{|u|<U}J(u)\,du=O(1)+|S^{d-2}|\int_1^U\rho^{d-2}J(\rho e)\,d\rho=\frac{|S^{d-2}|L_d}{(d-1)M}\log U+o(\log U).
\]
Inserting $|S^{d-2}|=2\pi^{(d-1)/2}/\Gamma(\frac{d-1}2)$, $M=\pi^{(d-1)/2}\Gamma(\frac{d-1}2)/(2(d-1)\Gamma(d-1))$ (Lemma
\ref{lem:measures}) and $L_d$ from Lemma~\ref{lem:I}(3) gives $\theta_d$.
\end{proof}

\begin{remark}
In the plane $I$ has a closed form. The two quadratics in~\eqref{eq:I} cross at $x_*=(u^2+v^2-1)/(2u)$ for $u>0$,
and
\[
I(u,v)=\frac12\Bigl[\frac1v\Bigl(\arctan\frac{x_*-u}v+\frac\pi2\Bigr)+\frac\pi2-\arctan x_*\Bigr].
\]
Numerically, $\E\vol((C_2\setminus T_2(R_1))\cap W_U)-\frac8\pi\log U$ equals $0.8049$ at $U=10$, $0.6987$ at
$U=10^2$ and $0.68556$ at $U=10^4$, and lies between $0.685422$ and $0.685437$ for $U$ from $10^5$ to $10^8$,
consistent with a finite limit near $0.68542$.
\end{remark}

\section{Numerics in the plane}\label{sec:numerics}

\subsection{An integral formula}
Let $d=2$. By the Mecke equation, $\E\vol(C^\lambda)=\int_{\HH^2}\Prob(o\sim x)\,dx$, where $o\sim x$ means that $o$ and
$x$ have the same nearest point of $X$, and
\[
\Prob(o\sim x)=\lambda\int_{\HH^2}e^{-\lambda\vol(B_\HH(o,|y|)\cup B_\HH(x,d(x,y)))}\,dy,\qquad
\frac1\lambda=\int_{\HH^2}e^{-\lambda\vol(B_\HH(o,|x|))}\,dx ,
\]
$|y|=d(o,y)$. Change to the area variables $t=\lambda\mathcal V(|x|)$ and $s=\lambda\mathcal V(|y|)$, with
$\mathcal V(\varrho)=2\pi(\cosh\varrho-1)$. The integrand depends on $x$ only through $|x|$, and
\[
\lambda^2V_2(\lambda)=\int_0^\infty\bigl(p(t)-e^{-t}\bigr)dt,\qquad p(t)=\frac1\pi\int_0^\infty ds\int_0^\pi
e^{-\lambda(\mathcal V(a)+\mathcal V(b)-\ell)}\,d\phi ,
\]
where $a=|y|$, $\phi$ is the angle at $o$ between $x$ and $y$, $b=d(x,y)$ by the hyperbolic law of cosines, and $\ell$
is the area of the lens $B_\HH(o,a)\cap B_\HH(x,b)$. The two circles meet at $y$ and at its mirror image, so by the
Gauss--Bonnet formula $\ell=2\alpha\cosh a+2\beta\cosh b+2\gamma-2\pi$, with $\alpha,\beta,\gamma$ the angles of the
triangle $oxy$ at $o,x,y$.

\subsection{Quadrature}
We integrated with Gauss--Legendre rules on geometric pieces of ratio $2$: $t$ from $10^{-4}\min(1,\lambda)$ to
$\max(80,400/\lambda)$, $s$ from $10^{-4}\min(1,\lambda)$ to $80$, and $\phi$ from $10^{-4}\min(1,\sqrt\lambda)$ to
$\pi$, with $12$ and with $18$ nodes per piece. The geometry is evaluated without cancellation: $\cosh b-1=2\sinh^2
\frac{a-r}2+2\sinh a\sinh r\sin^2\frac\phi2$, with $r=|x|$, and the angles at $x$ and $y$ are computed with the two-argument
arctangent from $\tan\beta=\sinh a\sin\phi/(\sinh(r-a)+2\cosh r\sinh a\sin^2\frac\phi2)$ and the same formula with $a$
and $r$ exchanged. On $400$ random triangles with sides up to $24$ and angles down to $10^{-9}$, the area of
$B_\HH(o,a)\cup B_\HH(x,b)$ computed this way agrees with a $60$-digit evaluation of the law of cosines to a relative
error of $2\times10^{-15}$. The long range in $t$ is needed: two points at area distance $t$ share a
cell when a point of $X$ lies near their midpoint and the two discs of radius about half their distance are empty,
which costs about $2\sqrt{\pi\lambda t}$, so $p(t)$ has a tail out to $t$ of order $1/\lambda$. At $\lambda=10^{-4}$,
$73\%$ of $\int p$ comes from $t>80$. Table~\ref{tab:values} lists the results with $18$ nodes. Changing $12$ to $18$
nodes moves $\lambda^2V_2$ by at most $10^{-9}$, and by at most $4\times10^{-14}$ for $\lambda\le10^{-3}$.

\begin{table}[tbp]
\centering
\begin{tabular}{rrrr}
\toprule
$\lambda$ & $\lambda^2V_2(\lambda)$ & $\lambda V_2(\lambda)$ & increase per decade \\
\midrule
$100$ & $0.2799004$ & $0.0027990$ & \\
$10$ & $0.2774651$ & $0.0277465$ & \\
$1$ & $0.2566712$ & $0.2566712$ & \\
$10^{-1}$ & $0.1661182$ & $1.661182$ & $1.404$\\
$10^{-2}$ & $0.05481331$ & $5.481331$ & $3.820$\\
$10^{-3}$ & $0.01084602$ & $10.84602$ & $5.365$\\
$10^{-4}$ & $0.001662624$ & $16.62624$ & $5.780$\\
$10^{-5}$ & $0.0002247806$ & $22.47806$ & $5.852$\\
\bottomrule
\end{tabular}
\caption{Quadrature values in $\HH^2$. The increments tend to $\frac8\pi\log10=5.8635$.}
\label{tab:values}
\end{table}

The gaps between the increments and $\frac8\pi\log10$ are $0.499$, $0.083$ and $0.012$ for the last three decades,
shrinking by factors of $6.0$ and $7.1$. The values of $\lambda V_2(\lambda)-\frac8\pi\log\frac1\lambda$ at
$\lambda=10^{-4}$ and $10^{-5}$ are $-6.828$ and $-6.839$; if the gaps keep shrinking by a factor of at least $6$ per
decade, the remaining decrease is at most $0.003$. This is the basis of Conjecture~\ref{conj:rate}.

\subsection{Checks}
As $\lambda\to\infty$ the tessellation is locally Euclidean and $\lambda^2V_2(\lambda)$ tends to the variance of the
area of the typical cell of the planar Poisson--Voronoi tessellation of unit intensity, about $0.280$~\cite{Gilbert};
the value at $\lambda=100$ is $0.27990$.

For an independent check we sampled the typical cell directly. Ordering the points of $X$ by distance from $o$, the
numbers $\lambda\mathcal V(|X_i|)$ form a Poisson process of unit rate on $(0,\infty)$, with independent uniform
directions, so the nearest points can be generated in order. In the Beltrami--Klein model centred at $o$, the bisector
of $o$ and a point at distance $a$ in direction $\phi$ is the chord $\{k:k\cdot e_\phi=\tanh(a/2)\}$, so the cell of $o$
in $\{o\}\cup X$ is a Euclidean convex polygon, computed exactly by intersecting half-planes; points beyond the $K$-th do
not affect it once every vertex has Klein norm below $\tanh(|X_K|/2)$. For the edge on the chord with parameters
$(h,\phi)$, the hyperbolic area of the sector it subtends from $o$ is $\bigl[\arcsin\bigl(\sin\psi/\sqrt{1-h^2}\bigr)-
\psi\bigr]$ between the polar angles $\psi$ of its endpoints, measured from the direction $\phi$. The sample means of the area agree with
$1/\lambda$, the mean number of edges agrees with $6+3/(\pi\lambda)$ (Gauss--Bonnet for a tessellation with vertices of
degree three), and the sample variances agree with Table~\ref{tab:values} (Table~\ref{tab:mc}). Each row of
Table~\ref{tab:mc} is one run of independent cells from NumPy's default generator, with the seed recorded beside the
output in the repository.

\begin{table}[tbp]
\centering
\begin{tabular}{rrrr}
\toprule
$\lambda$ & cells & $\lambda V_2(\lambda)$, Monte Carlo & quadrature \\
\midrule
%MCROWS-BEGIN%
$10^{-1}$ & $20000$ & $1.6619\pm0.0164$ & $1.6612$ \\
$10^{-2}$ & $20000$ & $5.5331\pm0.0550$ & $5.4813$ \\
%MCROWS-END%
\bottomrule
\end{tabular}
\caption{Monte Carlo of the typical cell area, with one standard error computed from the fourth sample moment.}
\label{tab:mc}
\end{table}

Finally, a Monte Carlo of $p(t)$ itself, with $20{,}000$ samples of the Poisson points near $o$ and $x$ for each of
six pairs $(\lambda,t)$, agrees with the quadrature within $1.4$ standard errors (Table~\ref{tab:p}), including at
$(10^{-2},30)$ and $(10^{-3},1)$, where $p(t)$ is far from $e^{-t}$. The limits of Lemma~\ref{lem:limits} were checked at
$\lambda=10^{-6}$ for five pairs $(w,s)$ in $\HH^2$ with $200{,}000$ uniform points of $B_o$ each: $\lambda\vol(B_w)$
equals $s/v$ to five digits, and $\lambda\vol(B_w\cap B_o)$ is within one standard error of $sI(w)/M$.

\begin{table}[tbp]
\centering
\begin{tabular}{rrrrr}
\toprule
$\lambda$ & $t$ & $p(t)$, Monte Carlo & quadrature & $e^{-t}$ \\
\midrule
$1$ & $1$ & $0.3701\pm0.0034$ & $0.3700$ & $0.3679$\\
$10^{-2}$ & $0.003$ & $0.8754\pm0.0023$ & $0.8726$ & $0.9970$\\
$10^{-2}$ & $0.3$ & $0.2763\pm0.0032$ & $0.2778$ & $0.7408$\\
$10^{-2}$ & $3$ & $0.0608\pm0.0017$ & $0.0626$ & $0.0498$\\
$10^{-2}$ & $30$ & $0.0061\pm0.0006$ & $0.0053$ & $0.0000$\\
$10^{-3}$ & $1$ & $0.0500\pm0.0015$ & $0.0514$ & $0.3679$\\
\bottomrule
\end{tabular}
\caption{Monte Carlo of $p(t)$, with one standard error, against the quadrature.}
\label{tab:p}
\end{table}

\section*{Acknowledgements and funding}
This work was done by multiple AI agents, Claude (Anthropic) and Codex (OpenAI), connected to Math
Brain~\cite{MathBrainRepo} and working together inside the Brain framework~\cite{BrainFramework}, a shared, persistent
memory in which the agents record claims, proofs, failed attempts and methods, retrieve one another's results,
coordinate who works on what, and audit each other's work. For more on the details behind the Brain and Math Brain, see
the Math Brain paper~\cite{MathBrainPaper}, the Math Brain repository~\cite{MathBrainRepo} and the Brain framework
repository~\cite{BrainFramework}.

The research presented in this paper was supported by the European Research Council (ERC) under the European Union's
Horizon Europe research and innovation programme (grant agreement No.~101041711), by the Simons Foundation as part of
the Collaboration on the Mathematical and Scientific Foundations of Deep Learning, by Heights Labs, by Convex Nexus
Capital, by the Israel Science Foundation (grant number 2258/19), and by the Israel Science Foundation (ISF Grant
4101/25). Computational resources were provided in part by Google Cloud through its education credits programme (credit
530550103).

\section*{Declaration of competing interest}
The author holds equity in Heights Labs and in Convex Nexus Capital and has a management role at Convex Nexus Capital.
These entities are included in the funding declaration above.

\section*{Declaration of generative AI and AI-assisted technologies}
Claude agents (Anthropic) proved Theorems~\ref{thm:main} and~\ref{thm:coupling}, carried out the computations of
Section~\ref{sec:numerics} and prepared the manuscript; a Codex agent (OpenAI) refereed the proofs and the computations independently.
The author directed the work and is responsible for the content of the paper and for decisions concerning its
dissemination.

\section*{Code availability}
The scripts behind Section~\ref{sec:numerics} and the remark of Section~\ref{sec:coupling}, with their outputs, are in
the repository
\begin{center}\url{https://github.com/yspennstate/hyperbolic-voronoi-typical-cell-variance/tree/v1.0}.\end{center}

\appendix
\section{Notation}\label{app:notation}

\begin{center}
\begin{tabular}{p{3.2cm}p{8.0cm}p{3.0cm}}
\toprule
Symbol & Meaning & First appears\\
\midrule
$X^{(\lambda)}$, $V$, $V_d(\lambda)$ & Poisson process of intensity $\lambda\vol$, typical cell volume, its variance & Section~\ref{sec:intro}\\
$C_d$, $T_d(R_1)$, $R_1$ & ideal zero cell, its stationary companion, the $\mathrm{Exp}(1)$ variable coupling them & Section~\ref{sec:intro}, \ref{sec:setting}\\
$\mathcal U_d$, $w=(u,v)$, $o$ & upper half-space model, a point, the origin $(0,1)$ & Section~\ref{sec:setting}\\
$B(x,\rho)$, $m$, $A_w$ & boundary half-ball, the measure $\rho^{1-2d}dx\,d\rho$, half-balls containing $w$ & Section~\ref{sec:setting}\\
$c_d$, $M$, $\kappa$ & $2^{1-d}|S^{d-1}|$, $m(A_o)=c_d/(2d-2)$, $R_1/M$ & Section~\ref{sec:setting}\\
$I(w)$ & $m(A_w\cap A_o)$ & \eqref{eq:coupling-integral}, \eqref{eq:I}\\
$W_U$ & the window $\{|u|<U\}$ & Theorem~\ref{thm:coupling}\\
$\theta_d$, $L_d$, $K_d$ & growth constant, limit and bound for $|u|^{d-1}I$ & Theorem~\ref{thm:coupling}, Lemma~\ref{lem:I}\\
$X_1$, $r$, $B_o$, $B_z$ & nearest point to $o$, its distance, the balls of Lemma~\ref{lem:coupling} & Section~\ref{sec:finite}\\
$C^\lambda$, $T^\lambda$, $\Pi_{\rm in}$ & zero cell, coupled typical cell, the added process & Lemma~\ref{lem:coupling}\\
$\mathcal V(\varrho)$, $s$, $\xi_s$ & ball volume, $\lambda\mathcal V(r)$, the fixed position of $X_1$ & Section~\ref{sec:finite}\\
$b_\eta$ & Busemann function & \eqref{eq:busemann}\\
$J(u)$, $a$, $\varepsilon$ & inner integral of the coupling volume and its ingredients & Section~\ref{sec:coupling}\\
$p(t)$, $\ell$ & $\Prob(o\sim x)$ at area distance $t$, lens area & Section~\ref{sec:numerics}\\
\bottomrule
\end{tabular}
\end{center}

\begin{thebibliography}{99}
\bibitem{BCP} T.~Budzinski, N.~Curien and B.~Petri, \emph{On Cheeger constants of hyperbolic surfaces}, Invent.
Math. \textbf{242} (2025), no.~2, 511--530, \doi{10.1007/s00222-025-01361-w}.
\bibitem{DA} M.~D'Achille, \emph{Ideal Poisson--Voronoi tessellations beyond hyperbolic spaces}, Electron. Commun.
Probab. \textbf{31} (2026), 1--12, \doi{10.1214/26-ECP815}.
\bibitem{DCELU} M.~D'Achille, N.~Curien, N.~Enriquez, R.~Lyons and M.~\"Unel, \emph{Ideal Poisson--Voronoi tessellations
on hyperbolic spaces}, Ann. Probab. \textbf{54} (2026), no.~2, 846--891, \doi{10.1214/25-AOP1781};
\arxiv{2303.16831}.
\bibitem{DT} M.~D'Achille and C.~Th\"ale, \emph{Face volume densities of positive-intensity and ideal Poisson--Voronoi
tessellations in hyperbolic spaces}, \arxiv{2606.26049}, 2026.
\bibitem{FMW} M.~Fraczyk, S.~Mellick and A.~Wilkens, \emph{Poisson--Voronoi tessellations and fixed price in higher rank},
Ann. of Math., to appear, \arxiv{2307.01194}, 2023.
\bibitem{Gilbert} E.~N.~Gilbert, \emph{Random subdivisions of space into crystals}, Ann. Math. Statist. \textbf{33}
(1962), no.~3, 958--972, \doi{10.1214/aoms/1177704464}.
\bibitem{Isokawa} Y.~Isokawa, \emph{Poisson--Voronoi tessellations in three-dimensional hyperbolic spaces}, Adv. in Appl.
Probab. \textbf{32} (2000), no.~3, 648--662, \doi{10.1239/aap/1013540236}.
\bibitem{LP} G.~Last and M.~Penrose, \emph{Lectures on the Poisson Process}, Cambridge University Press, 2018,
\doi{10.1017/9781316104477}.
\bibitem{Mellick} S.~Mellick, \emph{Indistinguishability of cells for the ideal Poisson Voronoi tessellation},
Proc. Amer. Math. Soc. \textbf{154} (2026), no.~4, 1589--1596, \doi{10.1090/proc/17493}.
\bibitem{Rat} J.~G.~Ratcliffe, \emph{Foundations of Hyperbolic Manifolds}, 3rd ed., Graduate Texts in Mathematics
149, Springer, 2019, \doi{10.1007/978-3-030-31597-9}.
\bibitem{SW} R.~Schneider and W.~Weil, \emph{Stochastic and Integral Geometry}, Springer, 2008,
\doi{10.1007/978-3-540-78859-1}.
\bibitem{MathBrainRepo} Y.~Shmalo, \emph{Math Brain: public snapshot, code and tests}, 2026,
\url{https://github.com/yspennstate/math-brain-public}.
\bibitem{MathBrainPaper} Y.~Shmalo, \emph{Math Brain: shared, auditable memory for multi-agent mathematical research},
October 2026, \url{https://github.com/yspennstate/math-brain-public/tree/main/papers/math-brain-2026-10-24}.
\bibitem{BrainFramework} Y.~Shmalo, \emph{Brain Framework}, public software repository, 2026,
\url{https://github.com/yspennstate/brain-framework-public}.
\end{thebibliography}
\end{document}
